% ============================================================================
% PAPER B · A Z[i]-module structure on body-centered-cubic corners,
%           with applications to two and only two of the order-4
%           occurrences in cubic-lattice algebra
% ============================================================================
% Target venue: Letters in Mathematical Physics (~10 pages).
% Strategy:     companion to PAPER_A_PI_FREE_GENERATOR.tex
% Created:      2026-05-02 afternoon (post FTD-0122 closure).
% Citation targets:
%   - docs/theory/09_mathematical/DERIV_BCC_COMPLEX_STRUCTURE.md (canonical)
%   - scripts/proofs/proof_bcc_complex_structure.py (verification)
%   - dissemination/papers/PAPER_A_PI_FREE_GENERATOR.tex (companion)
% Build:   cd dissemination/papers && pdflatex PAPER_B_BCC_COMPLEX_STRUCTURE.tex
% Status:  DRAFT v1 (2026-05-02).
%
% Scope (strict):
%   - Theorem 1: 8 BCC corners under 90-deg rotation -> 2 orbits of size 4.
%   - Theorem 2: Z[BCC] ⊗ Q  =  V_triv² ⊕ V_sign² ⊕ V_complex²;
%                V_complex carries natural Z[i]-module structure ≅ Z[i]².
%   - Theorem 3: Z[i]-module structure unifies Role 1 (CM Aut count =
%                |Z[i]^×| = 4) and Role 3 (master-quadratic tower level
%                k = 4, prefactor 16 = |Z[i]^×|²).
%   - Theorem 4 (no-go): Roles 2 (|O_h^ab| = 4 with O_h^ab = Klein four)
%                and 4 (27-block has 4 O_h orbits of different sizes)
%                CANNOT unify with Z[i]^× via group-theoretic homomorphism.
%
% Excluded (per scope):
%   - Master quadratic discovery / proof (deferred to Paper A)
%   - Number-field-theoretic results about G* (Paper A)
%   - Empirical α / N_c match (Paper A)
%   - Engine-level / EFT-program work
%   - Physical interpretation of 'why FTD prefers Z[i]'
% ============================================================================

\documentclass[11pt,reqno]{amsart}

\usepackage[margin=1in]{geometry}
\usepackage{amsmath,amssymb,amsthm,mathtools}
\usepackage{booktabs}
\usepackage{hyperref}
\usepackage{xcolor}
\usepackage[numbers,sort&compress]{natbib}
\usepackage{enumitem}

\hypersetup{colorlinks=true,
  linkcolor=blue!50!black,
  citecolor=blue!50!black,
  urlcolor=blue!50!black}

\theoremstyle{plain}
\newtheorem{theorem}{Theorem}[section]
\newtheorem{lemma}[theorem]{Lemma}
\newtheorem{proposition}[theorem]{Proposition}
\newtheorem{corollary}[theorem]{Corollary}
\theoremstyle{definition}
\newtheorem{definition}[theorem]{Definition}
\theoremstyle{remark}
\newtheorem{remark}[theorem]{Remark}

\numberwithin{equation}{section}

\newcommand{\Gstar}{G_{\!*}}
\newcommand{\Aut}{\operatorname{Aut}}
\newcommand{\End}{\operatorname{End}}
\newcommand{\disc}{\operatorname{disc}}
\newcommand{\rank}{\operatorname{rank}}
\newcommand{\BCC}{\mathrm{BCC}}
\newcommand{\Vtriv}{V_{\mathrm{triv}}}
\newcommand{\Vsign}{V_{\mathrm{sign}}}
\newcommand{\Vcpx}{V_{\mathrm{cpx}}}

\title[A {$\mathbb{Z}[i]$}-module structure on BCC corners]
{A $\mathbb{Z}[i]$-module structure on body-centered-cubic\\
 corners, with applications to two of the\\
 order-4 occurrences in cubic-lattice algebra}

\author{William J.\ cpaci-tani III}
\email{cpacitanimulli@gmail.com}

\author{C.\ Pacitanimulli}
\email{cpacitanimulli@gmail.com}

\date{\today}

\subjclass[2020]{Primary 20C10, 20C30;
                 Secondary 11G15, 11R52, 81V05.}
\keywords{Cubic lattice; permutation modules; Gaussian integers;
  complex multiplication; octahedral group; abelianization;
  representation theory of finite cyclic groups.}

\begin{document}

\begin{abstract}
Let $\BCC := \{(s_1, s_2, s_3) : s_i \in \{+1, -1\}\}$ denote the eight
corners of the unit cube and $J$ denote ninety-degree rotation in the
$(x,y)$ plane. Then $\langle J \rangle \cong \mathbb{Z}/4$ acts on
$\BCC$ with two orbits of size four; the integer permutation module
$\mathbb{Z}[\BCC] \cong \mathbb{Z}^{8}$ decomposes as
$\Vtriv^{2} \oplus \Vsign^{2} \oplus \Vcpx^{2}$ over $\mathbb{Q}$,
where the complex isotypic component $\Vcpx^{2}$ carries a natural
$\mathbb{Z}[i]$-module structure with $i$ acting as $J$, and is
isomorphic to $\mathbb{Z}[i]^{2}$ as a free $\mathbb{Z}[i]$-module of
rank two. This unifies two of the well-known order-four occurrences
in the algebra of the cubic lattice: the cardinality
$|\Aut_{\overline{\mathbb{Q}}}(E)| = |\mathbb{Z}[i]^{\times}| = 4$
of the automorphism group of the elliptic curve
$E\colon y^{2} = x^{3} - x$ (which has complex multiplication by
$\mathbb{Z}[i]$), and the level $k = 4$ at which a one-parameter
polynomial tower built from the lemniscatic constant first acquires
prefactor $2^{k} = 16 = |\mathbb{Z}[i]^{\times}|^{2}$. We further
record an explicit obstruction to extending the unification: the
abelianization $O_{h}^{\mathrm{ab}}$ of the full octahedral group is
isomorphic to the Klein four group $\mathbb{Z}/2 \times \mathbb{Z}/2$
rather than to $\mathbb{Z}/4$, so any homomorphism
$\mathbb{Z}[i]^{\times} \to O_{h}^{\mathrm{ab}}$ must factor through
the quotient $\mathbb{Z}[i]^{\times} / \{\pm 1\} \cong \mathbb{Z}/2$,
killing the order-four element. The $4$-element $O_{h}$-orbit
decomposition of the $27$-site $3 \times 3 \times 3$ block is
similarly an obstruction at the orbit-count level: the four orbits
have sizes $(1, 6, 12, 8)$, and no nontrivial finite group action
permutes them. Both occurrences are therefore order-or-count
coincidences with $|\mathbb{Z}[i]^{\times}|$ rather than
group-theoretic identifications. The paper records the partial
unification together with its precise no-go.
\end{abstract}

\maketitle

% ============================================================================
\section{Introduction}
\label{sec:intro}

The integer four appears in several distinct algebraic contexts
attached to the cubic lattice. We catalogue four such occurrences,
all known classically:

\begin{enumerate}[label=\textup{(\Alph*)}, leftmargin=*]
\item \emph{Elliptic-curve automorphisms.} The curve
$E\colon y^{2} = x^{3} - x$ has complex multiplication by the
Gaussian integers $\mathbb{Z}[i]$, and
$\Aut_{\overline{\mathbb{Q}}}(E) = \mathbb{Z}[i]^{\times}$ acts via
$(x, y) \mapsto (u^{2} x, u^{3} y)$ for $u \in \mathbb{Z}[i]^{\times}$;
hence $|\Aut_{\overline{\mathbb{Q}}}(E)| =
|\mathbb{Z}[i]^{\times}| = 4$.
\item \emph{Octahedral abelianization.} The full octahedral group
$O_{h}$ of order $48$ has commutator subgroup of index four; the
abelianization $O_{h}^{\mathrm{ab}}$ is therefore a group of order
four, and is isomorphic to the Klein four-group
$\mathbb{Z}/2 \times \mathbb{Z}/2$.
\item \emph{Lemniscatic polynomial tower.} For the lemniscatic
constant $\Gstar := \Gamma(1/4)/\Gamma(3/4)$ and integer $k \geq 3$,
the polynomial
$M_{k}(x) := x^{2} - 2^{k} \Gstar^{k-2} x + 2^{k} \Gstar^{k-1}$
has roots $x_{\pm}^{(k)}$ satisfying the harmonic invariant
$1/y_{+} + 1/y_{-} = 1$ for $y := x/\Gstar$, at every level $k$. The
level $k = 4$ is distinguished by being the smallest level whose
discriminant correction is irrational (conditional on Schneider's
1949 transcendence theorem) and at which the prefactor $2^{k} = 16$
matches $|\mathbb{Z}[i]^{\times}|^{2}$.
\item \emph{27-block orbit count.} The $3 \times 3 \times 3$ block
of sites in $\mathbb{Z}^{3}$ centred at the origin decomposes under
the natural $O_{h}$-action into exactly four orbits, of sizes one
(the centre), six (face neighbours), twelve (edge neighbours), and
eight (vertex neighbours, i.e.\ the BCC corners of the unit cube
inscribed in the block). The decomposition implies that the
multiplicity of the trivial irreducible representation $A_{1g}$ of
$O_{h}$ in the natural permutation representation on the 27 sites is
four.
\end{enumerate}

The four occurrences (A)--(D) are individually classical and
unrelated. The structural question we address is: \emph{do any of
them share a common group-theoretic origin}?

The main result of this paper is a partial unification together with
an explicit no-go for the rest. Specifically:

\begin{itemize}[leftmargin=*]
\item Occurrences (A) and (C) share the group-theoretic content
$\mathbb{Z}[i]^{\times}$, both descending from a natural
$\mathbb{Z}[i]$-module structure on a sublattice of the integer
permutation module of the BCC corners (Theorem~\ref{thm:bcc-cpx}).
\item Occurrences (B) and (D) do not. The abelianization
$O_{h}^{\mathrm{ab}}$ is the Klein four group, not the cyclic group
$\mathbb{Z}/4$, and the four 27-block $O_{h}$-orbits have
distinct sizes; both facts forbid a group-theoretic identification
with $|\mathbb{Z}[i]^{\times}| = 4$ (Theorem~\ref{thm:no-go}).
\end{itemize}

The paper is organised as follows. \S\ref{sec:bcc-action} fixes
notation for the BCC corners and the cyclic action of
ninety-degree rotation. \S\ref{sec:complex-structure} states and
proves the BCC complex-structure theorem (Theorem~\ref{thm:bcc-cpx}).
\S\ref{sec:cm} gives the connection to (A) via the
$\mathbb{Z}[i]$ endomorphism ring of $E$; \S\ref{sec:tower} gives the
connection to (C) via the lemniscatic tower of \cite{paper-a}.
\S\ref{sec:no-go} proves the no-go for (B) and (D).
\S\ref{sec:discussion} records the open questions the no-go does
\emph{not} settle.

\paragraph{Notation.}
We write $\mathbb{Z}[i]$ for the Gaussian integers, $\mathbb{Z}[i]^{\times} =
\{1, -1, i, -i\}$ for its unit group, $O_{h}$ for the full
octahedral group of order $48$ acting on $\mathbb{R}^{3}$ by signed
coordinate permutations, and $A_{1g}, A_{2g}, A_{1u}, A_{2u}, E_{g},
T_{2g}, T_{1u}, T_{2u}$ for its irreducible characters in standard
spectroscopic notation. For a finite group $G$ acting on a set $X$
we write $\mathbb{Z}[X]$ for the integer permutation module, viewed
as a $\mathbb{Z}[G]$-module.

% ============================================================================
\section{The BCC corners and the rotation action}
\label{sec:bcc-action}

\begin{definition}[BCC corners]
\label{def:bcc}
Let
\[
  \BCC \;:=\; \{(s_1, s_2, s_3) \,:\, s_i \in \{+1, -1\}\} \;\subset\; \mathbb{R}^{3}.
\]
Equivalently, $\BCC$ is the orbit of $(1, 1, 1)$ under the sign
group $\{\pm 1\}^{3}$. We have $|\BCC| = 8$.
\end{definition}

\begin{definition}[Rotation action]
\label{def:J}
Define $J \colon \mathbb{R}^{3} \to \mathbb{R}^{3}$ by
\[
  J(s_1, s_2, s_3) \;:=\; (-s_2,\; s_1,\; s_3).
\]
Then $J$ is rotation by ninety degrees in the $(x, y)$ plane and
preserves $\BCC$ setwise. The cyclic group $\langle J \rangle$ has
order four, and $J^{2}$ acts on $\BCC$ as the diagonal sign reversal
$(s_1, s_2, s_3) \mapsto (-s_1, -s_2, s_3)$.
\end{definition}

\begin{lemma}[Two orbits of size four]
\label{lem:orbits}
The cyclic group $\langle J \rangle$ acts on $\BCC$ with exactly two
orbits, each of size four:
\begin{align*}
  \mathcal{O}_{+} \;&=\; \{(+1,+1,+1),\, (-1,+1,+1),\, (-1,-1,+1),\, (+1,-1,+1)\}, \\
  \mathcal{O}_{-} \;&=\; \{(+1,+1,-1),\, (-1,+1,-1),\, (-1,-1,-1),\, (+1,-1,-1)\}.
\end{align*}
\end{lemma}

\begin{proof}
By Definition~\ref{def:J}, $J$ acts trivially on the third
coordinate, so the level sets $\{s_{3} = +1\}$ and $\{s_{3} = -1\}$
are invariant. On each such level set, the pair $(s_1, s_2)$ ranges
over four sign-combinations, on which $J$ acts as the four-cycle
\[
  (+,+) \;\to\; (-,+) \;\to\; (-,-) \;\to\; (+,-) \;\to\; (+,+).
\]
Each level set is therefore a single $\langle J \rangle$-orbit of
size four.
\end{proof}

\begin{remark}[Choice of plane]
The choice of the $(x, y)$ plane in Definition~\ref{def:J} is
conventional; rotation by ninety degrees in the $(y, z)$ or $(x, z)$
plane gives the same statement of Lemma~\ref{lem:orbits} with the
roles of the coordinates permuted. The constructions of
\S\ref{sec:complex-structure} are stated for the $(x, y)$ choice
without loss of generality.
\end{remark}

% ============================================================================
\section{The BCC complex-structure theorem}
\label{sec:complex-structure}

The integer permutation module $\mathbb{Z}[\BCC]$ is a free
$\mathbb{Z}$-module of rank eight on the basis $\{\delta_{v} : v \in
\BCC\}$, with $\langle J \rangle \cong \mathbb{Z}/4$ acting by
permutation: $J \cdot \delta_{v} = \delta_{J(v)}$.

\begin{theorem}[BCC complex structure]
\label{thm:bcc-cpx}
The $\mathbb{Q}\langle J \rangle$-module
$\mathbb{Z}[\BCC] \otimes_{\mathbb{Z}} \mathbb{Q}$ decomposes as a
direct sum of three isotypic components,
\begin{equation}
\label{eq:isotypic}
  \mathbb{Z}[\BCC] \otimes_{\mathbb{Z}} \mathbb{Q}
   \;=\;
  \Vtriv \;\oplus\; \Vsign \;\oplus\; \Vcpx,
\end{equation}
with $\mathbb{Q}$-dimensions $\dim \Vtriv = 2$, $\dim \Vsign = 2$,
$\dim \Vcpx = 4$. The component $\Vcpx$ carries a natural
$\mathbb{Z}[i]$-module structure with $i$ acting as $J$, and is
isomorphic, as a $\mathbb{Z}[i]$-module, to $\mathbb{Z}[i]^{2}$ of
rank two. Equivalently, $J|_{\Vcpx}^{2} = -\mathrm{id}|_{\Vcpx}$,
making $\Vcpx$ into a $\mathbb{Q}(i)$-vector space of dimension two.
\end{theorem}

\begin{proof}
We decompose orbit-by-orbit and assemble.

\medskip

\textit{Step 1: per-orbit decomposition.}
Let $\mathcal{O}$ be one of the two orbits of
Lemma~\ref{lem:orbits}, say $\mathcal{O} = \mathcal{O}_{+}$, and
write its elements in cyclic order $v_1, v_2, v_3, v_4$ with
$J v_{i} = v_{i+1 \!\bmod 4}$. The induced $\mathbb{Q}\langle J
\rangle$-module $\mathbb{Q}[\mathcal{O}] := \mathbb{Z}[\mathcal{O}]
\otimes \mathbb{Q}$ is the regular representation of
$\langle J \rangle \cong \mathbb{Z}/4$. The irreducible
representations of $\mathbb{Z}/4$ over $\mathbb{Q}$ are
\begin{itemize}[leftmargin=*]
\item the one-dimensional trivial representation (with $J \mapsto 1$),
\item the one-dimensional sign representation (with $J \mapsto -1$),
\item the two-dimensional ``rotation'' representation, which is
$\mathbb{Q}$-irreducible with minimal polynomial $T^{2} + 1$ but
splits over $\mathbb{Q}(i)$ into one-dimensional representations
with $J \mapsto +i, -i$.
\end{itemize}
Their $\mathbb{Q}$-dimensions sum to $1 + 1 + 2 = 4$, matching
$|\mathcal{O}|$, and a direct check verifies the regular
representation of $\mathbb{Z}/4$ over $\mathbb{Q}$ contains each
irreducible exactly once.

The corresponding orthogonal idempotents are
\begin{align*}
  \pi_{\mathrm{triv}} \;&=\; \tfrac{1}{4}\bigl(1 + J + J^{2} + J^{3}\bigr), \\
  \pi_{\mathrm{sign}} \;&=\; \tfrac{1}{4}\bigl(1 - J + J^{2} - J^{3}\bigr), \\
  \pi_{\mathrm{cpx}}  \;&=\; \tfrac{1}{2}\bigl(1 - J^{2}\bigr).
\end{align*}
One verifies $\pi_{\mathrm{triv}} + \pi_{\mathrm{sign}} +
\pi_{\mathrm{cpx}} = 1$, $\pi_{*} \pi_{\dagger} = 0$ for $* \neq
\dagger$, and $\pi_{*}^{2} = \pi_{*}$. The traces are $1, 1, 2$, so
the $\mathbb{Q}$-dimensions of the corresponding subspaces of
$\mathbb{Q}[\mathcal{O}]$ are $1, 1, 2$.

The trivial component is spanned by $v_{1} + v_{2} + v_{3} + v_{4}$;
the sign component by $v_{1} - v_{2} + v_{3} - v_{4}$. The complex
component is spanned by $\{a, b\}$ where
\[
  a \;:=\; v_{1} - v_{3}, \qquad b \;:=\; v_{2} - v_{4},
\]
on which $J$ acts as $J a = J(v_{1} - v_{3}) = v_{2} - v_{4} = b$
and $J b = J(v_{2} - v_{4}) = v_{3} - v_{1} = -a$. In the basis
$(a, b)$ the matrix of $J$ is
$\bigl(\begin{smallmatrix} 0 & -1 \\ 1 & \phantom{-}0 \end{smallmatrix}\bigr)$,
hence $J^{2} = -I_{2}$.

\medskip

\textit{Step 2: assembly over both orbits.}
Both orbits give isomorphic decompositions; summing over them yields
\eqref{eq:isotypic} with $\mathbb{Q}$-dimensions
$(2, 2, 4)$, summing to $|\BCC| = 8$.

\medskip

\textit{Step 3: $\mathbb{Z}[i]$-module structure on $\Vcpx$.}
Since $J^{2} = -\mathrm{id}$ on $\Vcpx$, the assignment
$i \mapsto J|_{\Vcpx}$ extends to a $\mathbb{Q}$-algebra homomorphism
$\mathbb{Q}(i) \to \End_{\mathbb{Q}}(\Vcpx)$, making $\Vcpx$ into a
$\mathbb{Q}(i)$-vector space. Its $\mathbb{Q}(i)$-dimension equals
$(\dim_{\mathbb{Q}} \Vcpx)/2 = 2$. The integer lattice
\[
  \Vcpx^{\mathbb{Z}} \;:=\;
  \pi_{\mathrm{cpx}} \cdot \mathbb{Z}[\BCC]
\]
is then a $\mathbb{Z}[i]$-stable lattice in a two-dimensional
$\mathbb{Q}(i)$-vector space, hence is isomorphic to a
$\mathbb{Z}[i]$-module of $\mathbb{Z}[i]$-rank two. Such modules are
classified by the structure theorem for finitely generated modules
over the principal ideal domain $\mathbb{Z}[i]$; absence of torsion
(any element with finite $\mathbb{Z}[i]$-annihilator would have
finite annihilator over the subring $\mathbb{Z} \subset
\mathbb{Z}[i]$, but $\Vcpx^{\mathbb{Z}}$ is torsion-free as a
$\mathbb{Z}$-submodule of $\mathbb{Z}[\BCC]$) and rank two over
$\mathbb{Z}[i]$ together force
$\Vcpx^{\mathbb{Z}} \cong \mathbb{Z}[i]^{2}$.
\end{proof}

\begin{remark}[On the absence of an ambient {$\mathbb{Z}[i]$}-structure]
\label{rem:no-zi-on-z3}
Theorem~\ref{thm:bcc-cpx} produces a $\mathbb{Z}[i]$-structure on a
sublattice of the eight-dimensional permutation module $\mathbb{Z}
[\BCC]$, not on the ambient three-dimensional integer lattice
$\mathbb{Z}^{3}$ in which the BCC corners sit. There is no integer
matrix $J \in M_{3}(\mathbb{Z})$ with $J^{2} = -I_{3}$: the minimal
polynomial of any such $J$ would be $T^{2} + 1$, and the
characteristic polynomial of a $3 \times 3$ matrix has degree three,
so the decomposition $\chi_{J}(T) = (T^{2} + 1)(T - \lambda)$ would
require a real eigenvalue $\lambda$ with $\lambda^{2} = -1$, which
is impossible. The $\mathbb{Z}[i]$-structure of
Theorem~\ref{thm:bcc-cpx} is therefore not a complex structure on
ambient three-dimensional space, but a $\mathbb{Z}[i]$-structure on
a four-dimensional subrepresentation of $\mathbb{Z}[\BCC] \otimes
\mathbb{Q}$. This subtlety is essential to the main result: the
naive ambient interpretation fails, and the per-orbit
representation-theoretic interpretation succeeds.
\end{remark}

\begin{remark}[Computational verification]
\label{rem:verification}
The orbit decomposition (Lemma~\ref{lem:orbits}), the per-orbit
projector relations, and the matrix identity $J|_{\Vcpx}^{2} = -I$
are verified in exact rational arithmetic in the script
\texttt{proof\_bcc\_complex\_structure.py} of the project
repository, available alongside this paper.
\end{remark}

% ============================================================================
\section{Connection (A): elliptic-curve automorphisms}
\label{sec:cm}

The first concrete order-four occurrence we identify with
Theorem~\ref{thm:bcc-cpx} is classical.

\begin{theorem}[CM connection]
\label{thm:cm}
Let $E$ be the elliptic curve $y^{2} = x^{3} - x$ over
$\overline{\mathbb{Q}}$. Then $E$ has complex multiplication by
$\mathbb{Z}[i]$, and the automorphism group
$\Aut_{\overline{\mathbb{Q}}}(E) \cong \mathbb{Z}[i]^{\times}$
acts via
$(x, y) \mapsto (u^{2} x,\, u^{3} y)$ for $u \in \mathbb{Z}[i]^{\times}$.
The four units $\{1, -1, i, -i\}$ produce four distinct $(u^{2},
u^{3})$ scaling pairs and hence
$|\Aut_{\overline{\mathbb{Q}}}(E)| = |\mathbb{Z}[i]^{\times}| = 4$.
This is the same group $\mathbb{Z}[i]^{\times}$ whose left-action
furnishes the $\mathbb{Z}[i]$-module structure on $\Vcpx$ in
Theorem~\ref{thm:bcc-cpx}.
\end{theorem}

\begin{proof}
The endomorphism ring computation $\End(E) = \mathbb{Z}[i]$ is
classical for the curve $y^{2} = x^{3} - x$ and is recorded in any
standard text on complex multiplication; see for example
\cite[Chapter II, \S 2]{silverman-cm}. The unit-action formula
$(x, y) \mapsto (u^{2} x, u^{3} y)$ is the $\mathbb{Z}[i]^{\times}$
sub-case of the standard formula for the action of an isogeny by
multiplication; the four units of $\mathbb{Z}[i]$ produce four
distinct scaling pairs because $u \mapsto u^{2}$ is injective on
$\mathbb{Z}[i]^{\times}$. The identification with $\mathbb{Z}[i]
^{\times}$ acting on $\Vcpx$ is then a tautology: both groups are
literally $\mathbb{Z}[i]^{\times}$, acting by left multiplication on
the corresponding $\mathbb{Z}[i]$-modules ($H^{0}(E, \Omega^{1})$ for
$E$, $\Vcpx^{\mathbb{Z}}$ for the BCC corners).
\end{proof}

\begin{remark}[On what Theorem~\ref{thm:cm} does not claim]
The connection of Theorem~\ref{thm:cm} is at the level of the
abstract group $\mathbb{Z}[i]^{\times}$, not at the level of any
specific geometric or arithmetic correspondence between $E$ and
$\BCC$. We do not claim any natural map between $E$ and $\BCC$, nor
do we claim that the automorphisms of $E$ are realised geometrically
on $\BCC$. The shared content is the unit group of the Gaussian
integers and its left action on $\mathbb{Z}[i]$-modules.
\end{remark}

% ============================================================================
\section{Connection (C): the lemniscatic polynomial tower}
\label{sec:tower}

The second order-four occurrence we identify is internal to the
algebra of $\Gstar = \Gamma(1/4)/\Gamma(3/4)$, in the sense of the
companion paper \cite{paper-a}.

\begin{definition}[Tower]
\label{def:tower}
For each integer $k \geq 3$, define
\begin{equation}
\label{eq:tower-def}
  M_{k}(x) \;:=\; x^{2} - 2^{k} \Gstar^{k-2} x + 2^{k} \Gstar^{k-1}.
\end{equation}
\end{definition}

\begin{theorem}[Tower-level connection]
\label{thm:tower-connection}
The level $k = 4$ is distinguished within the tower
\eqref{eq:tower-def} by two independent properties:
\begin{enumerate}[label=\textup{(\roman*)}, leftmargin=*]
\item The prefactor $2^{k} = 16 = 4^{2} =
|\mathbb{Z}[i]^{\times}|^{2}$, where $\mathbb{Z}[i]^{\times}$ is the
unit group acting on $\Vcpx$ in Theorem~\ref{thm:bcc-cpx}; for
$k \neq 4$ the prefactor $2^{k}$ is not the square of
$|\mathbb{Z}[i]^{\times}|$.
\item The discriminant correction $A_{k} := 2^{k-2} \Gstar^{k-3} -
1$ is rational for $k = 3$ and is conditionally transcendental over
$\mathbb{Q}$ for every $k \geq 4$ (conditional on the
Schneider-Chudnovsky theorem on the transcendence of $\Gamma(1/4)$),
making $k = 4$ the smallest level at which the discriminant
correction is irrational. See \cite[Theorem 3.2]{paper-a}.
\end{enumerate}
The two properties coincide at $k = 4$.
\end{theorem}

\begin{proof}
Property (i) is a calculation: $2^{k} = (2^{k/2})^{2}$, so $2^{k} =
n^{2}$ for $n \in \mathbb{Z}_{>0}$ iff $k$ is even, in which case
$n = 2^{k/2}$. The condition $n = |\mathbb{Z}[i]^{\times}| = 4$
forces $2^{k/2} = 4$, i.e.\ $k/2 = 2$, i.e.\ $k = 4$. Hence $k = 4$
is the unique level at which $2^{k}$ equals
$|\mathbb{Z}[i]^{\times}|^{2}$. Property (ii) is recorded as
\cite[Theorem 3.2]{paper-a} and the references therein; we do not
re-prove it here.
\end{proof}

\begin{remark}[Linking the two properties]
The two properties of Theorem~\ref{thm:tower-connection} are
\emph{a priori} unrelated: (i) is an elementary integer
computation, (ii) requires a theorem of contemporary transcendence
theory. They both happen to single out $k = 4$, but for
disconnected reasons in the present paper. We do not prove a
common origin; we record the coincidence as motivation for further
work.
\end{remark}

% ============================================================================
\section{No-go for occurrences (B) and (D)}
\label{sec:no-go}

We now show that the two remaining occurrences of the integer four
in \S\ref{sec:intro}--namely $|O_{h}^{\mathrm{ab}}| = 4$ and the
27-block orbit count--cannot share group-theoretic origin with
$|\mathbb{Z}[i]^{\times}| = 4$.

\begin{theorem}[Group-order obstruction at $O_{h}^{\mathrm{ab}}$]
\label{thm:no-go}
Let $\varphi \colon \mathbb{Z}[i]^{\times} \to O_{h}^{\mathrm{ab}}$
be any group homomorphism. Then $\varphi$ kills the elements of
order four in $\mathbb{Z}[i]^{\times}$; equivalently, $\varphi$
factors through the quotient $\mathbb{Z}[i]^{\times}/\{\pm 1\}
\cong \mathbb{Z}/2$. In particular, $\varphi$ is not injective. Any
identification of the order-four occurrence in
$O_{h}^{\mathrm{ab}}$ with that in $\mathbb{Z}[i]^{\times}$
therefore cannot be group-theoretic.
\end{theorem}

\begin{proof}
The unit group $\mathbb{Z}[i]^{\times} = \{1, -1, i, -i\}$ is
cyclic of order four, generated by $i$ with $i^{4} = 1$. The
abelianization $O_{h}^{\mathrm{ab}}$ has order four
\cite[\S\,VI.2]{johnston-group-theory} and is isomorphic to the
Klein four-group $\mathbb{Z}/2 \times \mathbb{Z}/2$, generated by
the cosets of inversion (parity) and a reflection (sign of
permutation), each of order two; hence every element of
$O_{h}^{\mathrm{ab}}$ has order one or two.

Let $\varphi$ be any group homomorphism. Since $i \in \mathbb{Z}[i]
^{\times}$ has order four, $\varphi(i) \in O_{h}^{\mathrm{ab}}$
satisfies $\varphi(i)^{4} = \varphi(i^{4}) = 1$, but the order of
$\varphi(i)$ divides both four (the order of $i$) and the order of
its image group; since every element of $O_{h}^{\mathrm{ab}}$ has
order one or two, $\varphi(i)$ has order one or two. If
$\varphi(i)$ had order four, then $\varphi(i)^{2}$ would have order
two, contradicting our characterisation of
$O_{h}^{\mathrm{ab}}$.

Hence $\varphi(i)$ has order one or two, equivalently $\varphi(i^{2})
= \varphi(i)^{2} = 1$, equivalently $-1 \in \ker \varphi$. The
kernel of $\varphi$ contains $\{\pm 1\}$, the unique subgroup of
$\mathbb{Z}[i]^{\times}$ of order two; the induced map $\overline{\varphi}
\colon \mathbb{Z}[i]^{\times} / \{\pm 1\} \to O_{h}^{\mathrm{ab}}$
is well-defined. Since $\mathbb{Z}[i]^{\times} / \{\pm 1\} \cong
\mathbb{Z}/2$ has order two and $|O_{h}^{\mathrm{ab}}| = 4$, the
map $\overline{\varphi}$ is at most a $2$-to-$1$ injection into a
size-four target; in particular $\varphi$ itself is not injective.
\end{proof}

\begin{theorem}[No-go at the 27-block orbit count]
\label{thm:no-go-orbits}
Let $G$ act on the four-element set $\Omega = \{O_1, O_2, O_3,
O_4\}$ of $O_{h}$-orbits in the 27-site block, where $|O_{1}| = 1$,
$|O_{2}| = 6$, $|O_{3}| = 12$, $|O_{4}| = 8$. Then $G$ acts
trivially on $\Omega$.
\end{theorem}

\begin{proof}
Any group element $g \in G$ inducing a nontrivial permutation of
$\Omega$ must permute orbits of equal $O_{h}$-size, since orbit
sizes are $G$-equivariant invariants of orbits. The four
$O_{h}$-orbits have pairwise distinct sizes $\{1, 6, 12, 8\}$, so
no two of them have equal size. Hence $g$ fixes each $O_{i}$ as a
set, and acts trivially on the four-element set $\Omega$. Since
the $G$-action on $\Omega$ is trivial, the count $|\Omega| = 4$ is
not a manifestation of any nontrivial group action; it is a
combinatorial count of $O_{h}$-orbit types in the 27-block.
\end{proof}

\begin{remark}[Net effect on the dual-four catalogue]
\label{rem:catalogue-net}
Theorems~\ref{thm:cm}, \ref{thm:tower-connection}, \ref{thm:no-go},
and \ref{thm:no-go-orbits} together establish:
\begin{itemize}[leftmargin=*]
\item Occurrences (A) and (C) are unified by the
$\mathbb{Z}[i]$-module structure of Theorem~\ref{thm:bcc-cpx}: in
both cases the order-four object is the unit group $\mathbb{Z}[i]
^{\times}$ acting by left multiplication on a $\mathbb{Z}[i]$-module
of small rank.
\item Occurrences (B) and (D) cannot be unified with (A), (C) at the
group-theoretic level: (B) is obstructed by the Klein-versus-cyclic
group-structure mismatch (Theorem~\ref{thm:no-go}); (D) is
obstructed by the orbit-size-mismatch combinatorial obstacle
(Theorem~\ref{thm:no-go-orbits}).
\end{itemize}
The four order-four occurrences therefore split as $2 + 2$: a
genuine pair sharing $\mathbb{Z}[i]^{\times}$ and a coincidental
pair contributing the integer $4$ at the level of group-order or
orbit-count alone.
\end{remark}

% ============================================================================
\section{Discussion}
\label{sec:discussion}

The partial unification of \S\ref{sec:complex-structure}--%
\S\ref{sec:tower} together with the no-go of \S\ref{sec:no-go}
leaves three open questions which we record without attempting to
resolve.

\paragraph{(1) Geometric origin of the $\mathbb{Z}[i]$-structure.}
Theorem~\ref{thm:bcc-cpx} produces a $\mathbb{Z}[i]$-module
structure on a sublattice of $\mathbb{Z}[\BCC]$, but the
construction depends on the choice of rotation plane (here the
$(x, y)$ plane). Rotation by ninety degrees in the $(y, z)$ or
$(x, z)$ plane yields an isomorphic statement with a different
$\mathbb{Z}[i]$-structure on a different sublattice. We do not
identify a canonical $\mathbb{Z}[i]$-structure on $\mathbb{Z}[\BCC]$
that is independent of the choice of axis, nor do we know whether
such a structure exists. (Note that any candidate must commute
with all three independent ninety-degree rotations, which are
elements of $O$ pairwise generating the rotation subgroup; the
constraint is nontrivial.)

\paragraph{(2) Maximality of the $\mathbb{Z}[i]$-sublattice.}
The complex-isotypic sublattice
$\Vcpx^{\mathbb{Z}} \subset \mathbb{Z}[\BCC]$ has $\mathbb{Z}$-rank
four and $\mathbb{Z}[i]$-rank two. We do not address whether
$\Vcpx^{\mathbb{Z}}$ is the unique $\mathbb{Z}[i]$-stable sublattice
of $\mathbb{Z}[\BCC]$ of these ranks; nor whether the
$\mathbb{Z}[i]$-isomorphism class of $\Vcpx^{\mathbb{Z}}$ is
canonical (Theorem~\ref{thm:bcc-cpx} establishes only that it is
free of rank two, and any free $\mathbb{Z}[i]$-module of given rank
is canonically isomorphic to $\mathbb{Z}[i]^{r}$).

\paragraph{(3) Higher tower levels.}
The tower of Definition~\ref{def:tower} has the harmonic-invariant
property at every level $k \geq 3$ \cite[Theorem 3.1]{paper-a}, but
only the level $k = 4$ has prefactor $|\mathbb{Z}[i]^{\times}|^{2}$
(Theorem~\ref{thm:tower-connection}(i)). Higher levels have
prefactors $2^{k}$ with $k \geq 5$, none of which is the square of
$|\mathbb{Z}[i]^{\times}|$. Whether any higher level admits a
distinct group-theoretic interpretation in $\mathbb{Z}[i]$ or in a
larger ring of integers (e.g.\ in the Eisenstein integers
$\mathbb{Z}[\omega]$, $\omega = e^{2\pi i / 3}$, with $|\mathbb{Z}
[\omega]^{\times}| = 6$ playing the role $|\mathbb{Z}[i]^{\times}|
= 4$ played here) is open. The companion paper \cite[\S 6.2]{paper-a}
reports a structural-uniqueness scan in which the Eisenstein-integer
multiplier family produces zero analogues of the dual-matching
property, suggesting that whatever distinguishes
$\mathbb{Z}[i]$ at level $k = 4$ is structurally specific to
$\mathbb{Z}[i]$ rather than generic to small rings of integers in
imaginary quadratic fields.

\paragraph{(4) Eisenstein-family null result and its interpretation.}
The Eisenstein-family null reported in \cite[\S 6.2]{paper-a} admits
a candidate but not a proof of structural origin: the orders of the
unit groups are $|\mathbb{Z}[i]^{\times}| = 4$ and $|\mathbb{Z}
[\omega]^{\times}| = 6$, and the four-versus-six distinction echoes
the Klein-versus-cyclic distinction
of Theorem~\ref{thm:no-go}: $\mathbb{Z}[\omega]^{\times} \cong
\mathbb{Z}/6$ is cyclic, not Klein, but it has order six and
contains an element of order three that has no analogue in
$\mathbb{Z}[i]^{\times}$. Whether this disparity in
unit-group structure is the structural reason for the Eisenstein
null, or whether the explanation lies elsewhere (e.g.\ in
discriminant or in the choice of CM curve), is open.

% ============================================================================
\section*{Acknowledgements}

The proof of Theorem~\ref{thm:bcc-cpx} reduces, after a brief
representation-theoretic step, to the elementary fact that the
regular representation of the cyclic group of order four over the
rationals decomposes as the direct sum of the trivial, the sign,
and a two-dimensional ``rotation'' irreducible. We use this fact
freely. Computational verification of the orbit decomposition, the
projector relations, and the matrix identity
$J|_{\Vcpx}^{2} = -I$ was carried out in exact rational arithmetic
in Python (\texttt{fractions.Fraction}); the verification script
\texttt{proof\_bcc\_complex\_structure.py} is available at the
project repository.

% ============================================================================
\begin{thebibliography}{99}

\bibitem{johnston-group-theory}
B.\ L.\ Johnston, \textit{Group representation theory and finite
groups}, lecture notes; standard reference for octahedral character
tables, e.g.\ in \cite[Chapter 9]{cornwell}.

\bibitem{cornwell}
J.\ F.\ Cornwell, \textit{Group Theory in Physics}, Volume 1: An
Introduction, Academic Press, 1997.

\bibitem{paper-a}
W.\ J.\ cpaci-tani III and C.\ Pacitanimulli, \textit{A $\pi$-free
generator of the lemniscatic field, with applications to a
closed-form quadratic for the fine-structure constant}, companion
paper, this issue (or arXiv preprint).

\bibitem{silverman-cm}
J.\ H.\ Silverman, \textit{Advanced Topics in the Arithmetic of
Elliptic Curves}, Graduate Texts in Mathematics \textbf{151},
Springer-Verlag, New York, 1994.

\end{thebibliography}

\end{document}
